% ML4H 2026, proceedings track (8 pages).
\documentclass[pmlr,twocolumn,10pt]{jmlr}

\mlhtrack{proceedings}

\newif\iffinal
\finalfalse  % \finaltrue for camera-ready

\iffinal
    \ifmlhneedspmlr
      \jmlrvolume{XXX}
      \jmlryear{2026}
    \fi
    \ifmlhfindings \jmlrproceedings{}{ML4H 2026 - Findings Track}\fi
    \ifmlhdemo     \jmlrproceedings{}{ML4H 2026 - Demo Track}\fi
    \jmlrworkshop{Machine Learning for Health (ML4H) 2026}
\else
    \jmlrproceedings{}{Submitted to ML4H 2026: \mlhtrackname}
    \jmlrworkshop{Machine Learning for Health (ML4H) 2026}
\fi

\urlstyle{same}

\usepackage{booktabs}
\usepackage{siunitx}
\usepackage[switch]{lineno}

\theorembodyfont{\upshape}
\theoremheaderfont{\scshape}
\theorempostheader{:}
\theoremsep{\newline}
\newtheorem*{note}{Note}

\title[Auditing the Auditor]{Auditing the Auditor: Four Failure Modes of
Confound Detection in Clinical Machine Learning}

\author{Anonymous Author(s)}

\linenumbers

\begin{document}

\maketitle

\ifmlhdemo\else
\begin{abstract}
Clinical machine learning results are routinely invalidated by confounds that
leave no trace in the reported metrics: a label definition that shifted between
sites, a pretraining corpus that overlapped the target, a hyperparameter swept
on out-of-distribution data. We built five automated detectors, one per confound
type, validated primarily against ground truth from real historical defects
rather than synthetic injection, with controlled injection reserved for the two
confounds that are inherently dose-response quantities. On a
demonstration-scale internal suite drawn from MIMIC-IV, eICU-CRD and the
PhysioNet 2019 corpus the detectors reach precision 1.00 and a false-positive
rate of 0.00 at recall 0.86.

That headline is not the contribution. Running the detectors against themselves
produced a second, more durable result: a taxonomy of four ways a diagnostic
fails \emph{independently of whether the confound it targets is real}. Each type
has a measured instance here. A diagnostic with no undecidable state reported
false cleanliness four times in independent parts of our own code, once
manufacturing a perfect confusion matrix from an undefined metric. A
conventional agreement threshold turned out to sit on the population value it
was meant to separate from, flagging 48.4\% of the time at full scale -- worse
than at 1/640th the data, not better. A detector that catches a real historical
bug in our code reached a verdict on zero of 727 third-party files. And the same
undecidable-state pattern is structurally present in DomainBed's model-selection
code, in a file that guards against it elsewhere.

We argue the failure modes generalize past clinical data and past our detectors,
and that the reason is specific: the pattern needs a component that renders a
verdict, and evaluation code mostly renders numbers.
\end{abstract}
\begin{keywords}
confound detection, reproducibility, clinical machine learning, label shift,
data leakage, model selection, negative results
\end{keywords}
\fi

\ifmlhneedsstatements
\paragraph*{Data and Code Availability}
We use three credentialed critical-care datasets from
PhysioNet~\citep{goldberger2000physiobank}: the PhysioNet 2019 Sepsis Challenge
corpus~\citep{reyna2020early}, MIMIC-IV~v3.1~\citep{johnson2023mimiciv} and
eICU-CRD~v2.0~\citep{pollard2018eicu}, each available to researchers who
complete the required training and sign the data use agreement. Our code is
available in an anonymized repository, \url{https://anonymous.4open.science/r/reboot-8568/}. It
contains all detector code and tests, the pre-registration and outcome
documents, run logs, and every aggregate result (confusion matrices, $\kappa$
values, confidence intervals, $t$-statistics); author-identifying commit history
is removed for review. Per-stay derived caches are not released because they are
individual-encounter artifacts of restricted data; the repository documents how
to regenerate them under the reader's own credentials, and the reproduction path
does not depend on them.

\paragraph*{Institutional Review Board (IRB)}
This work uses only de-identified, publicly available datasets obtained under
their respective data use agreements, and involves no interaction with human
subjects and no access to identifiable private information. It therefore does
not constitute human subjects research and did not require IRB approval.
% TODO: confirm IRB wording with institution.
\fi

\section{Introduction}
\label{sec:intro}

A clinical machine learning result can be wrong in a way that nothing in the
paper reveals. The model is sound, the code runs, the numbers are real, and the
finding is an artifact of how the data was assembled. Systematic reviews of
COVID-19 prediction models found essentially none fit for clinical
use~\citep{wynants2020prediction,roberts2021common}, and the failures were
overwhelmingly of this kind: not modeling error, but confounded construction.
Leakage in particular has been argued to underlie a reproducibility crisis
spanning many scientific fields that use machine
learning~\citep{kapoor2023leakage}.

The natural response is tooling. If a confound leaves a detectable signature,
write a detector for it, run it before publishing, and report the result. We did
this for five confound types we had personally encountered:

\begin{enumerate}
\itemsep0em
  \item \textbf{Label-definition shift} --- the operational definition of the
    outcome differs between sites.
  \item \textbf{Pretraining leakage} --- the pretraining corpus overlaps the
    evaluation site, so ``zero-shot'' is not.
  \item \textbf{OOD-contaminated selection} --- a hyperparameter is swept
    against the held-out distribution.
  \item \textbf{Circular derived constraints} --- a physiological constraint is
    imposed on a variable that was itself computed from that constraint.
  \item \textbf{Missingness and scale artifacts} --- a cross-site score
    difference is driven by which variables were recorded, not by physiology.
\end{enumerate}

Each detector is validated against a positive case \emph{and} a false-positive
control, with ground truth taken wherever possible from real historical defects
in our own repository rather than from synthetic injection. On the
demonstration-scale internal suite (\tableref{tab:main}) the detectors reach
precision 1.00 and FPR 0.00 at recall 0.86; \sectionref{sec:measured} shows a
full-scale control that flags on 48.4\% of resampled draws while its single
recorded verdict reads clean.

\paragraph{The contribution is what happened next.} Taking the detectors
seriously meant auditing them the way they audit a pipeline, and that audit
found more than the detectors did. A diagnostic can fail in ways that have
nothing to do with whether its target confound exists, and those failures are
invisible in exactly the metric a detector reports --- a clean confusion matrix.
We name four such failure modes (\sectionref{sec:taxonomy}), give a measured
instance of each (\sectionref{sec:measured}), and show that at least one
replicates outside clinical machine learning entirely
(\sectionref{sec:generalize}).

Two of these are uncomfortable results about our own instrument, and we report
them as results rather than as limitations. Detector~5 fails on its own positive
case. Detector~3 catches a real historical bug in our code and then reaches no
verdict at all on 727 files from 18 published third-party repositories. We
believe the second is the more useful finding of the two.

\section{Related Work}
\label{sec:related}

\citet{damour2022underspecification} show that a training pipeline can be
underspecified: many models fit the data equally well while behaving very
differently under distribution shift, so the reported validation number
underdetermines what was learned. Our concern is adjacent but distinct. We are
not asking whether the \emph{model} is underdetermined by the pipeline; we are
asking whether the \emph{diagnostic} is underdetermined by its own validation.
A detector reporting TN$=$2 tells you it did not fire twice. It does not tell
you whether it could have.

The closest structural analogue is
\citet{salaudeen2025misspecified}, who derive conditions under which a domain
generalization benchmark can detect what it claims to and then audit real
benchmarks against those conditions, finding that many are misspecified. Our
work sits one level further out: they ask whether a benchmark can measure what
it claims about a model, we ask whether a diagnostic can measure what it claims
about a pipeline. The two failure sets are related but not the same --- a
misspecified benchmark yields an uninformative number, whereas a vacuous
diagnostic yields an affirmatively \emph{wrong} one, because ``did not detect''
reads as ``nothing to detect''.

The confounds our detectors target are instances of shortcut
learning~\citep{geirhos2020shortcut} and of hidden
stratification~\citep{oakdenrayner2020hidden} in the specific sense that the
learned signal is real in the assembled data and absent in the deployment
setting. That literature characterises the phenomena; we are asking what happens
to the instruments built to find them.

Work on leakage~\citep{kapoor2023leakage} and on clinical model
appraisal~\citep{wynants2020prediction,roberts2021common} catalogues confounds
and recommends checklists. Checklists are read by humans and are not
falsifiable; a detector is executable and can be wrong in a way you can measure.
That is the gap this paper occupies, and the cost of closing it is that the
detector now needs its own validity argument --- which is where the four failure
modes come from. Among existing tooling, AIF360~\citep{bellamy2019aif360} is the
closest analogue in that it emits real verdicts rather than numbers for a human
to read, and \sectionref{sec:generalize} reports what we found when we looked
for our own failure modes inside it.

\section{A Taxonomy of Diagnostic Failure Modes}
\label{sec:taxonomy}

We define a \emph{diagnostic} as a procedure that audits an ML pipeline for one
specific confound and returns a verdict. The four types below are ways a
diagnostic fails regardless of whether its target confound is present.

\paragraph{Type 1 --- Vacuous pass.} A check that never reached a verdict is
scored as a verdict of ``clean''. The absence of a detection is recorded as a
detection of absence. Two sub-forms share this consequence.
\emph{Type 1a, falsy-as-clean}: an undecidable or absent result --- \texttt{nan},
\texttt{None}, an empty container --- is consumed by a boolean test, where it is
indistinguishable from ``no confound found''.
\emph{Type 1b, degenerate control}: the control cannot fail, passing for
arithmetic or ceiling reasons rather than because the detector discriminated,
and the resulting true negative is counted as evidence of specificity.

\paragraph{Type 2 --- Structural non-discriminability.} The decision boundary
coincides with, or sits arbitrarily close to, the true value it must separate
from. No amount of data resolves this, because it is not a sampling problem.

\paragraph{Type 3 --- Granularity or pairing mismatch.} The diagnostic is blind
to a confound whose structure does not match the diagnostic's aggregation level
or pairing structure.

\paragraph{Type 4 --- Post-hoc fitting risk.} Redesigning a diagnostic after
observing its failure is contaminated evidence, even when the redesign is well
reasoned. Only decision rules fixed in advance escape this.

Types 1 and 2 are failures of the verdict; types 3 and 4 are failures of the
design process that produced it. All four are compatible with a perfect
confusion matrix, which is why we report the matrix and the taxonomy together.

\section{Detectors and Validation Protocol}
\label{sec:methods}

\paragraph{Detectors.} Detector~1 computes Cohen's $\kappa$~\citep{cohen1960kappa}
on a fixed 500-patient audit subset between two operationalizations of the
outcome, and flags below the conventional Landis--Koch ``substantial agreement''
boundary of $0.60$~\citep{landis1977measurement}. The reference definition is
window-mode SOFA per Sepsis-3~\citep{singer2016sepsis3}, which specifies a SOFA
rise over an interval around the suspicion time rather than a reading at one
instant. Detector~2 trains a small masked-reconstruction probe on a source site
with a controlled fraction of target-site stays leaked in, and tests whether
probe loss on held-out target data falls; the test is a one-sided paired $t$
within seed with $df$ one less than the seed count (5 seeds, $df=4$). Detector~3 is static analysis: it traces, via
Python's abstract syntax tree, whether a sweep function's evaluation call
receives a loader derived from the held-out distribution. Detector~4 traces
variable lineage through the same AST machinery, with guard-condition reasoning,
to find constraints imposed on variables derived from those same constraints.
Detector~5 compares per-component availability ratios and a
composition-controlled aggregate across sites.

\paragraph{Ground truth.} Wherever possible the positive case is a real
historical defect. Detector~3's positive is a specific buggy sweep from our own
git history, retained verbatim alongside its corrected version. Detector~4's
positive is a genuine Severinghaus circularity in our PhysioNet loader, where
PaO\textsubscript{2} is derived from SaO\textsubscript{2} and then constrained
against it. Detector~1's positive is ICD-coded sepsis against SOFA-window
scoring on real cohorts. Only detector~2's leakage and detector~5's ablation are
injected, because the quantity being detected is a dose.

\paragraph{Controls are measurements, not identities.} Every detector ships a
false-positive control, and each control was checked for the degeneracy of
Type~1b. This check is not optional and it is not cosmetic: it invalidated our
first version of detector~1's control (\sectionref{sec:measured}).

\paragraph{Uncertainty is mandatory.} Every stochastic detector reports an
uncertainty measure with its verdict, or the table renders \texttt{MISSING}
rather than a blank. Detectors 3 and 4 are deterministic static analysis and are
excluded explicitly rather than silently. This is enforced in code: the harness
classifies each detector as stochastic or deterministic and refuses to classify
an unknown one.

\paragraph{Pre-registration.} Detector~5's design variant was decided by a rule
committed 29 minutes \emph{before} the commit adding its external-validation
code; the pre-registration file is unchanged since, and the outcome is recorded
in a separate document. Both commits are included as anonymized patches with
their full hashes (\appendixref{apd:prereg}).

\paragraph{Scale, stated per result.} We state the scale of every result
explicitly rather than declaring a default and marking exceptions, because an
unlabelled number sitting in an authoritative position is precisely the failure
this paper is about.

\tableref{tab:main} is the \emph{internal validation suite} and runs at
demonstration scale: detector~1 on the eICU demonstration cohort, detector~2 on
PhysioNet Site~A${\to}$Site~B at 900 stays per site, detector~5 on a
1200-stay-per-arm sample of the PhysioNet corpus. Detectors 3 and 4 are static
analysis over code and have no data scale at all; their scale axis is corpus
size (18 repositories, 727 parsed files for detector~3; nine constraint/loader
pairs for detector~4).

\emph{External} validation --- \sectionref{sec:results} onward --- runs at full
dataset scale on MIMIC-IV~v3.1, eICU-CRD~v2.0 and the complete PhysioNet 2019
corpus, and every such figure is marked full-scale where it appears. Model
capacity is held fixed across scales so the two remain comparable.

\section{Results}
\label{sec:results}

\begin{table}[htbp]
\floatconts
  {tab:main}
  {\caption{\textbf{Internal} validation suite, at demonstration scale (see
   \sectionref{sec:methods}). Each row is one detector against its positive case
   and its false-positive control(s). Full-scale external results are reported
   in the text and are \emph{not} folded into this matrix. Detector~5's row is a
   false negative, discussed in \sectionref{sec:measured}.}}
  {\footnotesize\setlength{\tabcolsep}{2.5pt}%
  \begin{tabular}{@{}lccl@{}}
  \toprule
  \bfseries Detector & \bfseries TP/FP/FN/TN & \bfseries Data & \bfseries Note\\
  \midrule
  1 label shift  & 1/0/0/1  & eICU demo & $\kappa$ 0.467/0.771\\
  2 leakage      & 3/0/0/1  & PhysioNet & floor 5\%\\
  3 OOD select.  & 1/0/0/1  & code      & hist.\ bug\\
  4 circularity  & 1/0/0/8  & code      & 9 pairs\\
  5 missingness  & 0/0/1/1  & PhysioNet & \bfseries false neg.\\
  \midrule
  \bfseries Total & \bfseries 6/0/1/12 &  & \\
  \bottomrule
  \end{tabular}}
\end{table}

\tableref{tab:main} reports TP/FP/FN/TN per detector for the internal suite.
The aggregate is $\mathrm{TP}=6$, $\mathrm{FP}=0$, $\mathrm{FN}=1$,
$\mathrm{TN}=12$: precision $1.00$, recall $0.86$, FPR $0.00$. Detector~1's
$\mathrm{TN}=1$ is the corrected control, SOFA window-mode against single-point
scoring at $\kappa=0.771$. The degenerate control of \sectionref{sec:measured}
(instance~1, $\kappa(x,x)=1$, $\mathrm{TN}=2$) was replaced before this matrix
was computed and appears in no table.

The external controls described below are deliberately \emph{not} added to that
matrix. Two of the four external controls per database are near-ceiling
(\sectionref{sec:measured}, instance~2), so folding them in would raise
detector~1's TN from 1 to 7 on the strength of tests that discriminate nothing
--- the Type~1b denominator inflation this paper argues against. We report them
in prose with their $\kappa$ values instead.

\paragraph{Detector 1 at full scale.} On full MIMIC-IV ($74{,}829$ stays) and
full eICU-CRD ($132{,}900$ stays), scoring a 500-patient audit draw, the
positive case and the discriminating control move as follows:

\begin{center}\footnotesize
\begin{tabular}{@{}lcc@{}}
\toprule
\bfseries Case & \bfseries demo & \bfseries full\\
\midrule
MIMIC ICD vs.\ SOFA-window (pos.) & 0.150 & 0.204\\
MIMIC window vs.\ single (ctrl)   & 0.651 & 0.634\\
eICU ICD vs.\ SOFA-window (pos.)  & 0.467 & 0.514\\
eICU window vs.\ single (ctrl)    & 0.771 & 0.869\\
\bottomrule
\end{tabular}\end{center}

\noindent Every verdict is unchanged: $\mathrm{TP}=1$, $\mathrm{FP}=0$,
$\mathrm{FN}=0$, $\mathrm{TN}=3$ on both databases at both scales.

A $5000$-resample bootstrap over the same cohorts gives \emph{whole-cohort}
$\kappa$, a different estimand: MIMIC positive $0.174$, control
$\mathbf{0.602}$; eICU positive $0.451$, control $0.826$. Each audit-draw value
lies inside its audit-subset bootstrap interval ($0.634 \in [0.539, 0.663]$;
$0.869 \in [0.737, 0.901]$), so the two agree but are not interchangeable;
operating-point figures here are always audit-subset.

\tableref{tab:uncertainty} gives the mandatory uncertainty column. Three
cautions attach to it and all three matter. Detector~1's interval is the
\emph{audit-subset} interval, because the audit subset is what the detector
actually scores; the whole-cohort interval is $[0.597, 0.607]$, ten times
tighter, and quoting it would make a coin-flip detector look decisive.
Detector~2 is reported as a paired $t$ rather than as a mean and standard
deviation because its unpaired null spread is comparable in size to the effect
it must reveal --- the unpaired form cannot see its own signal, which is itself
an instance of Type~3. Detector~5's flag rate is \emph{not} scale-invariant: the
same procedure gives 2/5 at $n=1200$ stays per arm and 0/3 at $n=4000$, so the
rate is a property of the audit sample size as much as of the sites and must
never be quoted without its $n$.

\begin{table}[htbp]
\floatconts
  {tab:uncertainty}
  {\caption{Uncertainty per stochastic detector; deterministic ones say why
   they have none. Det.~1: MIMIC-IV window vs.\ single-point control, full
   scale, 500-patient audit draws. Det.~2: MIMIC-IV$\to$eICU, full scale, 5\%
   leakage, 5 training seeds ($df=4$). Det.~5: PhysioNet Site~A vs.~B positive,
   demonstration scale; its seeds are random stay-sample draws, not training
   runs.}}
  {\small
  \begin{tabular}{@{}ll@{}}
  \toprule
  \bfseries Det. & \bfseries Uncertainty\\
  \midrule
  1 & 95\% CI $[0.539, 0.663]$, $P(\text{flag})=0.484$\\
  2 & paired $t=-6.65$ vs crit.\ $-2.132$, 5/5 signs\\
  3 & deterministic (no CI: static analysis)\\
  4 & deterministic (no CI: static analysis)\\
  5 & flagged 2/5 seeds at $n=1200$/arm\\
  \bottomrule
  \end{tabular}}
\end{table}

Detector~2's full-scale external run (MIMIC-IV source, $74{,}607$ stays; eICU
target, $130{,}446$ stays) detects at the lowest injected level, 5\% leakage,
with paired relative probe-loss deltas of $-3.2\%$, $-7.4\%$ and $-21.1\%$ at
5\%, 20\% and 100\% leakage and a clean 0\% control. Detector~4 returns
precision $1.00$ and recall $1.00$ over nine cases: the single genuine
Severinghaus circularity in the PhysioNet loader, and eight non-circular
constraint/loader pairs across all three databases.

\paragraph{Provenance caveat.} Detector~1's full-scale figures are transcribed
from run output; detectors 2 and 5 are verified against archived logs. The logs
for detector~1's full-scale run were lost in transfer. The numbers were read
directly off the run and the regeneration path is released and verified, but a
reader cannot check them against an artifact the way they can for the other two.
We state this rather than quietly presenting all three detectors on the same
evidential footing.

\section{The Failure Modes, Measured}
\label{sec:measured}

\subsection{Type 1: four vacuous passes in our own code}

A diagnostic with no undecidable state will eventually report false
cleanliness. When a check cannot reach a judgement --- the metric is undefined,
the pattern it looks for is absent, the control is too easy to fail --- that
outcome silently becomes ``no confound found'' and is then counted as evidence
of specificity. We hit this four times, independently, in our own code:

\begin{enumerate}
\itemsep0em
  \item \emph{(1b)} Detector~1's negative control passed the same audit array as
    both arguments to Cohen's $\kappa$. The result was $1.000$ by construction,
    and $\mathrm{TN}=2$ measured arithmetic rather than the detector.
  \item \emph{(1b)} Detector~1's suspicion-window width controls sit at $\kappa$
    $0.93$--$1.00$; reported bare, the external suite's $\mathrm{TN}=3$ (not
    \tableref{tab:main}) implies three tests when only one discriminates.
  \item \emph{(1a)} Detector~3's \textsc{indeterminate} verdict on third-party
    repositories. Folded into TN it would have produced $\mathrm{TN}=3$ from
    files the detector never analysed.
  \item \emph{(1a)} A baseline check's \texttt{train\_test\_gap}: one seed's
    in-domain AUROC was undefined, \texttt{nan > threshold} evaluated
    \texttt{False}, and the abstention was scored as a correct silence --- handing
    the baseline a perfect $\mathrm{TP}=3, \mathrm{FP}=0, \mathrm{FN}=0,
    \mathrm{TN}=1$ that would have contradicted this paper's own headline
    comparison.
\end{enumerate}

Four occurrences across independent code is not coincidence. The natural Python
idiom for ``did not detect'' and for ``could not tell'' is the same falsy value.
The fix is cheap and mechanical: every check returns a decidability flag
alongside its verdict, and undecidable cases are excluded from the confusion
matrix and reported in their own column rather than defaulting to either class.

A fifth instance appeared while preparing this manuscript. Two committed logs,
\texttt{logs/check5.out} and \texttt{logs/run\_all.out}, still reported
detector~5's composition figure as a $0.49$ ``share'' and closed with
\textsc{detector validated}, for the case this paper reports as its one false
negative. The $0.49$ had been retracted as an alphabetical-prefix sampling
artifact (population value ${\approx}0.27$, below the gate); the prose was
corrected, the machine output was not. Both logs are regenerated, and a test now
scans committed logs for the retracted wording. Correcting a claim while leaving
its evidence behind makes the least trustworthy statement in a repository the one
carrying the strongest provenance signal.

\subsection{Type 2: a threshold on the value it must separate}

Detector~1 flags when $\kappa$ on the audit subset falls below $0.60$, the
conventional Landis--Koch boundary, fixed before any of this was measured. On
full MIMIC-IV ($74{,}829$ stays) the \emph{whole-cohort} $\kappa$ between SOFA
window-mode and SOFA single-point scoring --- two equally defensible Sepsis-3
operationalizations --- is $\mathbf{0.602}$, 95\% CI $[0.597, 0.607]$. The population value sits on the
decision boundary. At the detector's operating point the control therefore flags
$\mathbf{48.4\%}$ of the time: the verdict is close to arbitrary.

This is not small-sample noise, and the small-sample data actively understated
it. At $n=117$ the same control gave $P(\text{flag})=0.216$; scaling the cohort
by $640\times$ more than \emph{doubled} the flag probability. More data narrows
the interval around $0.602$; it cannot move $0.602$ away from $0.60$.

We did not move the threshold. Adjusting it after observing this failure would
be exactly the Type~4 error our pre-registration exists to prevent, and the
temptation is stronger here precisely because the required adjustment is tiny.

The counterexample keeps this from reading as ``$\kappa$-based label-shift
detection does not work'': on the same run, eICU's equivalent control moved
\emph{away} from the boundary, $0.771 \rightarrow 0.869$ on the audit draw
(whole-cohort $0.826$; $P(\text{flag})=0.000$). The failure is database- and
axis-specific.

\subsection{Type 3: granularity and pairing}

Two of the five detectors initially produced false negatives, and both for the
same reason: the diagnostic did not mirror the structure of what it audits.
Detector~5 aggregated at the wrong granularity; detector~2 compared unpaired
where the comparison is paired. Detector~2's fix is why it is reported as a
paired $t$: its unpaired null spread is comparable to the effect size, so the
unpaired form could not see its own signal.

Full scale sharpened this rather than repeating it. Detector~5 gates on both an
availability ratio and a composition term. At full scale the composition gate
binds \emph{nowhere} --- every case clears its threshold --- so the conjunction
collapses into the availability condition alone and one of the two gates is
decorative. At demonstration scale the composition gate bound exactly once, so
the redundancy was invisible. A gate that does no work stays hidden until the
diagnostic is examined at the scale it will actually be used at.

Detector~5's own positive case remains a false negative: on PhysioNet
Site~A~vs.~Site~B the composition ratio is $0.27$ against a $0.30$ gate, just
under. On external data with controlled ablation it reaches precision $1.00$ and
FPR $0.00$ with a detection floor between 50\% and 80\% ablation across five
seeds. Both facts stand; neither cancels the other.

\begin{note}
The composition quantity is a \emph{ratio of two gaps}, not a share. It is
unbounded above~1 and exceeds $1.4$ on one control (a five-seed mean, PhysioNet
Site~A against itself). It must not be described as
``composition explains $X\%$ of the gap''.
\end{note}

\subsection{Type 4: what pre-registration bought}

We pre-registered a rule for choosing between two variants of detector~5 before
running either on external data. The alternative variant had higher recall at demonstration scale
($0.73$ vs $0.60$) and was tempting; the pre-registered criterion was a strictly
lower detection floor, which it did not deliver, so it was rejected. At full
scale the two variants are \emph{identical on every scored case}, and the
rejection that looked like a judgement call became a non-question.

The pre-registration also excluded one case from scoring in advance because its
ground truth is unknown. That case is the only one on which the two variants
diverge. Had it been scored, it alone would have decided the comparison --- on
data with no ground truth. Excluding it beforehand is the entire value of having
written the rule down first.

\section{Does This Generalize?}
\label{sec:generalize}

\subsection{Detector 3 reaches no verdict on 727 third-party files}

Detector~3 catches a real historical bug in our repository. To test whether that
is a detector or a project-specific pattern match, we scanned 18 public
repositories, each associated with a published paper, at pinned commits, cloned
read-only and parsed --- never executed.

Per-file ground truth requires a human reading each file and does not scale.
\emph{Recognition} does: the detector either identifies a function as a sweep it
can audit, or it does not, and that is mechanically decidable. So the scan
measures how often the detector reaches a judgement at all.

Of 735 Python files, 727 parsed and 8 failed to parse (reported in their own
column, never as clean). The detector recognised a sweep in
\textbf{zero} of the 727 --- both as written and after fitting its recognition
vocabulary to those very repositories, an extension from 19 names to 60
across five categories. The failure is structural, not lexical.
DomainBed~\citep{gulrajani2021search} illustrates why: it selects by
\texttt{argmax} over a table of already-recorded metrics, with no training call
anywhere in the selecting function, so there is no loop for the detector to
find.

The honest bound is therefore: detector~3 detects OOD-contaminated selection in
code shaped like ours, and does not reach a verdict on other codebases'
selection idioms at all.

\subsection{Type 1 outside clinical machine learning}

We read three non-clinical repositories at pinned commits for the
falsy-as-clean pattern: DomainBed and robustdg~\citep{mahajan2021domain} from the
18-repository corpus above, and AIF360~\citep{bellamy2019aif360}, added for this
search. None was executed.

DomainBed's query library maps \texttt{nan} to $-\infty$ inside its
\texttt{sorted} routine and ships a \texttt{filter\_not\_nan} helper, so the
hazard was anticipated in that file. Its \texttt{argmax}, used by three of the
four model-selection methods to pick the winning run, is a bare \texttt{max}
with no such guard; because every comparison against \texttt{nan} is false, a
\texttt{nan} score wins or loses purely on list position. Separately, its
leave-one-out selection guards missing environments with
\texttt{any([v == -1 for v in val\_accs])}, which is vacuously false on an empty
list --- the same shape as our own vacuous pass.

\textbf{Neither is reachable with shipped datasets}, since both require a
single-environment dataset and the minimum shipped is three. We report them as
structurally present and not currently exploitable. Claiming a live bug would be
the same overreach this paper exists to catch.

The value of the DomainBed case is precisely that the guard and its omission sit
in one file: the hazard was known, and one call site was still missed.

robustdg gives a variant. A domain with zero support is correctly guarded and
therefore never enters the metric dictionary; the aggregation then takes its key
set from run~0 alone. The same degeneracy is a silent gap if it hits the first
seed and a \texttt{KeyError} if it hits a later one. AIF360's empty-subgroup path
is deliberate and correct, and we record it as searched-and-not-found rather
than as a near miss.

\subsection{Why the pattern concentrates in diagnostics}

The most useful result of the cross-domain search is structural. A falsy value
can only be misread as ``clean'' where something asks a yes/no question.
Benchmarks overwhelmingly compute a number, format it, and print it for a human,
and a \texttt{nan} there is loud: it appears as \texttt{nan} in the results
table. Model selection is the one place benchmark code renders a verdict --- and
that is exactly where the guarding turns out to be inconsistent. Read this as an
argument for why diagnostic and selection code specifically needs an explicit
undecidable state; three repositories cannot establish how common the pattern is
in evaluation code generally.

\section{Comparison with Standard Practice}
\label{sec:baselines}

Would ordinary methodological hygiene have caught these confounds? We ran three
standard checks --- an external performance floor, $k$-fold instability, and the
train/test gap --- against detector~2's leakage injection, scored identically
(\tableref{tab:baselines}).

\textbf{This section is at demonstration scale} (2400 stays per site for the
baselines, 900 for detector~2). Its claim is a signal-to-nuisance ratio measured
on the same data, so it is internally consistent at this scale; whether it holds
at full scale is untested, and a full-scale rerun was out of budget. The
detector~2 deltas used here (PhysioNet A$\to$B: $-3.1/-9.0/-15.6\%$ at
5/20/100\% leakage) and the full-scale ones in \sectionref{sec:results}
(MIMIC$\to$eICU: $-3.2/-7.4/-21.1\%$) differ in database pair as well as scale,
so they are not a scale comparison.

\begin{table}[htbp]
\floatconts
  {tab:baselines}
  {\caption{Standard checks vs.\ a purpose-built detector on the same leakage
   injection, \textbf{demonstration scale} (2400 stays/site, 3 seeds). Recall
   $1.00$ here is \emph{not} competence.}}
  {\small
  \begin{tabular}{@{}lcc@{}}
  \toprule
  \bfseries Check & \bfseries TP/FP/FN/TN & \bfseries FPR\\
  \midrule
  external floor       & 3/1/0/0 & 1.00\\
  $k$-fold instability & 3/1/0/0 & 1.00\\
  train/test gap       & 3/1/0/0 & 1.00\\
  \midrule
  detector 2           & 3/0/0/1 & \bfseries 0.00\\
  \bottomrule
  \end{tabular}}
\end{table}

All three baselines fire at \emph{every} level including 0\% leakage. FPR is
$1.00$ and TN is $0$: a check that always fires has recall $1.00$ by
construction and carries no information.

The reason they are useless is quantitative, not directional. Leakage does move
downstream performance in the expected direction --- cross-site AUROC rises and
the in-domain/cross-site gap shrinks from $0.117$ to $0.103$ at this scale ---
but the ordinary
source-to-target domain gap at \emph{zero} leakage is $0.117$, an order of
magnitude larger than the ${\sim}0.01$ leakage-induced change. The confound is
buried under the baseline domain shift. Detector~2's probe loss on the same
data, by contrast, moves by the PhysioNet deltas above with a clean 0\%
control. The argument for purpose-built detectors is signal-to-nuisance ratio.

This comparison was run against detector~2's scenario only; detector~5's
confound is not a performance-degradation quantity and was out of scope.

\section{Discussion}
\label{sec:discussion}

\paragraph{A clean confusion matrix is not evidence that a detector works.} Every
failure in \sectionref{sec:measured} is compatible with the matrix in
\tableref{tab:main}, and three of them were present in versions of the code that
produced clean matrices. The minimum a diagnostic should report alongside its
verdict is: an explicit undecidable count, the discriminating power of each
control, an uncertainty measure at the operating point, and the decision rule's
provenance.

\paragraph{Scale changes diagnostics qualitatively, not just statistically.} Two
of our findings invert with data volume. Detector~1's near-miss got worse at
$640\times$ the cohort. Detector~5's composition gate went from occasionally
binding to never binding, collapsing a two-gate conjunction into a one-gate
test. Neither is visible at demonstration scale, and both change what the
detector is.

\section{Limitations}
\label{sec:limitations}

Types 2 and 3 each rest on a single instance, and we separate them by depth of
investigation rather than by count. Type~2's instance --- detector~1's
$\kappa=0.602$ against a $0.60$ boundary --- carries a $5000$-resample bootstrap,
a disconfirmed ``more data will resolve it'' hypothesis, and a cross-database
counterexample. Type~3's rests on two of our own initial false negatives plus
one full-scale observation about a redundant gate, with no external instance at
all, and is accordingly the type most likely to be revised. Neither should be
read as carrying the evidentiary weight of Type~1, which has four internal
occurrences and an external one.

Detector~3's generalization bound is a negative result, not a repaired detector.
We did not rewrite it to recognise DomainBed-style selection, because doing so
after observing the failure is the Type~4 error.

Detector~2 cannot separate ``the encoder saw target values'' from ``the encoder
saw target missingness pattern'', because the target site records more sparsely
and adding target stays moves both together. A disambiguating arm is designed
and costed but deliberately not run; it is a stated scope limit, not a gap we
were unaware of.

Detector~5's external component set is analogous, not identical: PhysioNet
compares SpO\textsubscript{2} with SaO\textsubscript{2}, while the external
databases lack SaO\textsubscript{2} and use the Severinghaus relation.

Detector~1's full-scale numbers are transcribed rather than archived
(\sectionref{sec:results}). Finally, all five detectors were built by the same
authors who then audited them; the external corpus scan and the cross-domain
read are our attempts to break that circularity, and they are partial.

\section{Conclusion}
\label{sec:conclusion}

We set out to build confound detectors for clinical machine learning and did:
five of them, precision $1.00$, FPR $0.00$ and recall $0.86$ on the
demonstration-scale internal suite, validated against real historical defects
and re-run at full dataset scale. Auditing those detectors the way
they audit a pipeline produced the more portable result --- four failure modes a
diagnostic exhibits independently of whether its target confound is real, each
with a measured instance, one of them structurally present in a widely used
non-clinical benchmark that guards against it elsewhere in the same file.

The practical recommendation is narrow and cheap: give every diagnostic an
explicit undecidable state, check that every control is a measurement rather
than an identity, report uncertainty at the operating point, and fix the
decision rule before looking. Each caught something in this project.

\acks{Acknowledgments go here \emph{but should only appear in the camera-ready
version of the paper if it is accepted}.}

\bibliography{ref}

\appendix

\section{Reproduction}\label{apd:repro}

The released repository reproduces \tableref{tab:main} from a clean clone with
no clinical data present: detectors 3 and 4 run live as static analysis, and
detectors 1, 2 and 5 load committed aggregate results. The accompanying test
suite runs 116 tests, of which 7 \emph{skip} with the explicit reason
\texttt{PhysioNet data not present} rather than passing. We note the distinction
because it is the same one this paper argues for: a test that cannot run has not
succeeded.

\section{Pre-registration}\label{apd:prereg}

The pre-registration and outcome documents are in the anonymized repository,
with anonymized patches of the two commits that fix their order
(\texttt{detectors/preregistration\_evidence/}): the commit adding the
pre-registration (\texttt{63f5d95a8a8e}, 2026-08-06 00:34) and its direct child
adding the variant rules and external-validation code (\texttt{1600d91071cf},
01:03). The repository's pre-registration file is identical to the one added by
the first commit. Author names are removed from the patches. Their dates are
author-supplied metadata and cannot be verified during review; the full hashes,
listed in that directory, are the check. A git hash covers a commit's author,
timestamp, content and parent, so the history released with the de-anonymized
repository must reproduce them. \sectionref{sec:measured} summarises the
outcome.


\end{document}
